% ch4_answers.tex -- Module 8, Chapter 4: Quick check answers.
\begin{chapteranswers}{4}{Properties of the Nucleus}
  % 4.1 Radioactive decay
  \item $Z = 53$ protons and $N = 131 - 53 = 78$ neutrons. In $\beta^-$ decay $A$ stays
        the same and $Z$ goes up by 1. $Z = 54$ is xenon:
        \final{$\chFourNuc{131}{53}{I} \to \chFourNuc{131}{54}{Xe} + \chFourBeta + \chFourAnu$}.
        Check: $131 = 131 + 0$ and $53 = 54 + (-1)$ \cmark.
  \item Work backwards. Before the $\beta^-$ decay, $A$ was the same (60) and $Z$ was
        one less ($28 - 1 = 27$): $\chFourNuc{60}{27}{Co}$. Before the neutron was
        absorbed, $A$ was one less ($60 - 1 = 59$) and $Z$ was the same.
        So \final{X $= \chFourNuc{59}{27}{Co}$ (cobalt-59)}.
  \item Subtract the background (30) from each reading: no absorber 420, paper 418,
        aluminium 15 counts per minute. Paper makes almost no difference, so there is
        \final{no alpha}. Aluminium removes most of the count, so there is
        \final{beta}. A small count (15 above background) still passes \qty{5}{mm} of
        aluminium, so there is also some \final{gamma}. The excess is small, so count
        for longer to confirm it. If it then falls only slowly with lead, it is gamma.
  % 4.2 Half-life and activity
  \item $\dfrac{960}{120} = 8 = 2^3$, so 3 half-lives passed in \qty{4.5}{h}.
        $t_{\frac{1}{2}} = \dfrac{\qty{4.5}{h}}{3}$, so
        \final{$t_{\frac{1}{2}} = \qty{1.5}{h}$}.
  \item $\lambda = \dfrac{\ln 2}{t_{\frac{1}{2}}} = \dfrac{0.6931}{\qty{3.82}{d}}
        = \qty{0.1814}{d^{-1}}$. Then
        $N_\mathrm{t}/N_0 = e^{-\lambda t} = e^{-0.1814 \times 10.0} = e^{-1.814}
        = 0.163$. Answer: \final{\qty{16.3}{\percent} is left}. Check: \qty{10.0}{d} is
        2.6 half-lives, and \qty{16}{\percent} lies between \qty{25}{\percent} and
        \qty{12.5}{\percent} \cmark.
  \item $\lambda = \dfrac{0.6931}{\qty{5730}{\yr}} = \qty{1.210e-4}{\yr^{-1}}$.
        Put $N_\mathrm{t} = 0.62\,N_0$ into the decay law and divide by $N_0$:
        $0.62 = e^{-\lambda t}$. Take ln: $\ln 0.62 = -0.4780 = -\lambda t$, so
        $t = 0.4780 \div \lambda = 0.4780 \div \num{1.210e-4} = \qty{3950}{\yr}$.
        Answer: \final{about \qty{4.0e3}{\yr} old} (2 significant figures, because
        \qty{62}{\percent} has 2). Check: more than half is left, so the age is less
        than one half-life (\qty{5730}{\yr}) \cmark.
  \item Each nucleus decays on its own, at random, with the same chance per second
        ($\lambda$). The other nuclei do not change this chance. So in any time
        interval the same \emph{fraction} of the nuclei decays, whether there are a
        thousand or $10^{20}$. The time for half of them to decay is therefore the same.
        A bigger sample has a bigger activity, but the same half-life.
  % 4.3 Mass defect, binding energy and energy released
  \item $Z = 3$, $N = 7 - 3 = 4$. Parts:
        $3 \times 1.007276 + 4 \times 1.008665 = \qty{7.056488}{\amu}$.
        $\Delta m = 7.056488 - 7.014356 = \qty{0.042132}{\amu}$.
        $E_\mathrm{B} = 0.042132 \times 931.5 = \qty{39.246}{MeV}$, and
        $E_\mathrm{B}/A = 39.246 \div 7 = \qty{5.607}{MeV}$.
        Answers: \final{$E_\mathrm{B} = \qty{39.2}{MeV}$},
        \final{\qty{5.61}{MeV} per nucleon}.
  \item Before: $14.003074 + 4.002603 = \qty{18.005677}{\amu}$. After:
        $16.999132 + 1.007825 = \qty{18.006957}{\amu}$. (Electrons: $7 + 2 = 8 + 1$, so
        atomic masses are fine.)
        $E = (18.005677 - 18.006957) \times 931.5 = -0.001280 \times 931.5$, so
        \final{$E = \qty{-1.19}{MeV}$}. The minus sign means the products have
        \emph{more} mass than the reactants. Energy must be supplied: the alpha
        particle must bring in more than \qty{1.19}{MeV} of kinetic energy. (In fact
        it needs at least about \qty{1.5}{MeV}. Momentum must be conserved, so the
        products must keep some kinetic energy.)
  \item In alpha decay there are only two products, the alpha particle and the
        daughter. The energy released is fixed, and conservation of momentum fixes how
        it is shared. So the alpha particles have fixed energies. (If the daughter is
        left in an excited state, the alpha has a little less energy and a gamma ray
        carries the rest.) In beta decay
        there are three products: the beta particle, the (anti)neutrino and the
        daughter. The fixed energy is shared between them in a different way in each
        decay. So the beta particle can have any energy from almost 0 up to
        $K_\mathrm{max}$.
  % 4.4 Fission and fusion
  \item Top numbers: $235 + 1 = 137 + 97 + x$, so $x = 236 - 234$:
        \final{$x = 2$}. Check the bottom numbers: $92 + 0 = 52 + 40$ \cmark.
  \item The \textbf{moderator} (for example water or graphite) slows the fast
        neutrons from fission. It does this by many collisions with light nuclei.
        Slow (thermal) neutrons are much more likely to be captured by U-235. Then it
        splits, and the chain reaction can continue. The \textbf{control rods} (boron
        or cadmium) absorb neutrons. Moving them in or out changes how many are
        absorbed. This keeps the average at exactly one neutron per fission causing a
        new fission, so the power stays steady. Pushed fully in, they stop the
        reactor.
  \item Before: $2.014102 + 3.016049 = \qty{5.030151}{\amu}$. After:
        $4.002603 + 1.008665 = \qty{5.011268}{\amu}$. (Electrons: $1 + 1 = 2$.)
        $E = (5.030151 - 5.011268) \times 931.5 = 0.018883 \times 931.5 = \qty{17.59}{MeV}$.
        Answer: \final{$E = \qty{17.6}{MeV}$}, the same as the binding-energy method in
        the formula card \cmark.
  \item Iron-56 sits at (very near) the peak of the curve, with about \qty{8.8}{MeV}
        per nucleon, one of the highest values of any nucleus. A nucleus with
        $A = 112$ lies to the right of the
        peak, with a lower value (about \qty{8.5}{MeV}). So the product would be
        \emph{less} tightly bound than the iron. It would have more mass than the two
        iron nuclei, so energy would have to be supplied, not released.
  \item (a) Energy of one fission:
        $\qty{200}{MeV} \times \qty{1.602e-13}{J.MeV^{-1}} = \qty{3.204e-11}{J}$.
        Fissions per second $= \dfrac{P}{E_\mathrm{fission}}
        = \dfrac{\qty{2.0e7}{W}}{\qty{3.204e-11}{J}} = \qty{6.242e17}{s^{-1}}$.
        Answer: \final{\num{6.2e17} fissions each second}.\newline
        (b) Fissions in a year:
        $\qty{6.242e17}{s^{-1}} \times \qty{3.156e7}{s} = \num{1.970e25}$.
        Mass: $\num{1.970e25} \times \qty{3.90e-25}{kg} = \qty{7.683}{kg}$.
        Answer: \final{about \qty{7.7}{kg}} (2 significant figures, like \qty{20}{MW}).
        Check: the rule ``\qty{1}{g} gives \qty{1}{MW} for a day'' gives
        $20 \times 365 \times \qty{1}{g} = \qty{7.3}{kg}$, close \cmark.
\end{chapteranswers}
